% GENERATED FILE. Edit canonical lessons, then run npm run generate:books.
% canonical-source-hash: fnv1a64:44a46feca59b14fa
% canonical-lessons: TA-C19-vayathu, TA-C19-age-register-grammar

\chapter{Asking Someone's Age}
\label{ch:ta-asking-age}

This chapter is generated from the canonical micro-lessons used by Language Ladder.
Each section stays independently resumable and preserves the authored prerequisite order.

\section[\ta{உனக்கு} \ta{எத்தனை} \ta{வயது} \ta{ஆகிறது}?]{\ta{உனக்கு} \ta{எத்தனை} \ta{வயது} \ta{ஆகிறது}? — Tamil actually has BOTH shapes}

\label{lesson:TA-C19-vayathu}



\begin{quote}
\textbf{Warm-up.} {[}PAUSE 2s{]} Same Sanskrit word one more time — but be careful here: Tamil isn't locked into just one of the two grammatical shapes you've now seen. It has both, in different registers.
\end{quote}

\subsection*{You'll want to know: \ta{வயது} — the same Sanskrit "vigor/age" word, a fourth time}
\textbf{\ta{வயது}} (\emph{vayathu}) = "\textbf{age}" — the same Sanskrit \textbf{\dv{वयस्}} (\emph{vayas}, "age, vigor," PIE cousin of Latin \textbf{vīs}) you've now met four times across Kannada, Telugu, Malayalam, and Tamil. All four Dravidian languages borrowed this exact Sanskrit word for "age" — none of them built a native alternative for this particular concept, unlike some others you've seen in this course.

\subsection*{Guided Practice}
{[}PAUSE 1s{]}

\begin{itemize}
\raggedright
  \item {[}YOU SAY: "vayathu" — age{]}
  \item {[}YOU CONNECT: Sanskrit \emph{vayas} · Tamil \emph{vayathu} · Latin \emph{vīs}{]}
\end{itemize}

\begin{tcolorbox}[breakable,colback=teal!4,colframe=teal!35!black,title={Wrap-up recall}]
{[}PAUSE 3s{]} Is Tamil's \textbf{\ta{வயது}} the same Sanskrit word borrowed by all four Dravidian languages here? (\textbf{Yes} — none built a native alternative for this concept.) What did \emph{vayas} originally include besides age? (\textbf{Vigor}.) Next: choose between Tamil's casual possessive and formal dative-plus-become shapes.
\end{tcolorbox}
\section[\ta{உன்} \ta{வயசு} \ta{என்ன}? / \ta{உனக்கு} \ta{எத்தனை} \ta{வயது} \ta{ஆகிறது}?]{Your age what? / To you how many years becomes?}

\label{lesson:TA-C19-age-register-grammar}



\begin{quote}
\textbf{Warm-up.} {[}PAUSE 2s{]} Tamil did not choose one age construction. Register selects between two shapes already visible elsewhere in the family.
\end{quote}

\begin{grammarlens}[title={Grammar Lens: Casual possessive}]
\begin{quote}
\textbf{\ta{உன்} \ta{வயசு} \ta{என்ன}?} — "your age what?"
\end{quote}

Casual speech uses a plain possessive with no dative and no verb, like Kannada and Telugu.
\end{grammarlens}

\begin{grammarlens}[title={Grammar Lens: Formal dative plus "become"}]
\begin{quote}
\textbf{\ta{உனக்கு} \ta{எத்தனை} \ta{வயது} \ta{ஆகிறது}?} — "to you how many years is becoming?"
\end{quote}

>

\begin{quote}
\textbf{\ta{எனக்கு} \ta{இருபது} \ta{வயது} \ta{ஆகிறது}.} — "to me twenty years is becoming."
\end{quote}

Formal/written Tamil uses the dative experiencer plus \textbf{\ta{ஆகிறது}} (\emph{aakirathu}, "is becoming," from \emph{aaku}). Malayalam uses a similar dative shape but chooses "exists" instead.

Do not compress these into "Tamil's one grammar." Both live in the language; setting and register choose between them.
\end{grammarlens}

\subsection*{Guided Practice}
{[}PAUSE 1s{]}

\begin{itemize}
\raggedright
  \item {[}YOU SAY: "un vayasu enna?" — casual{]}
  \item {[}YOU SAY: "unakku ettanai vayathu aakirathu?" — formal{]}
  \item {[}YOU ANSWER: "enakku irupathu vayathu aakirathu"{]}
\end{itemize}

\begin{tcolorbox}[breakable,colback=teal!4,colframe=teal!35!black,title={Wrap-up recall}]
{[}PAUSE 3s{]} Which shape is casual? (\textbf{Possessive}: "your age what?") Which is formal? (\textbf{Dative + aakirathu}, "become.") Does Malayalam use the same verb? (\textbf{No} — it uses "exists.")
\end{tcolorbox}
